\chapter{Pray}

“I don’t know,” Tibor McMasters said at last. He hesitated, weighing the words. Perhaps he had already gone too far. Too many attempts had been made on his life. Maybe reality itself was signaling him to turn back.

The bird waited.

“Let me tell you more,” it said. “Someone is still following you. His name is Pete.”

“Why?” Tibor asked. The knowledge stirred only a dull unease now. “What does he want?”

“I can’t determine that,” the bird replied. “But he means you no harm. Can you decide now?”

Tibor asked the question that mattered. “If I find the God of Wrath—will he try to kill me?”

“He will not know who you are at first,” the bird said. “Too much time has passed. He no longer believes anyone is pursuing him.”

That made sense. Tibor drew a slow breath. “Where is he?” he asked. “Lead me there. Slowly.”

“A hundred miles north,” the bird said. “You will find either him—or someone who resembles him.”

“That’s not very reassuring.”

“Nor is it certain,” the bird admitted.

Tibor felt disappointed. All this danger, and the reward was partial knowledge at best. But perhaps that was the nature of Pilg: fragments accumulating toward something barely coherent.

“What does he look like?” Tibor asked.

“Bad,” the bird said simply. “Old. Stooped. Fat. Missing teeth. Foul breath.”

The description stunned Tibor. The Deus Irae—reduced to a decaying man.

“There’s nothing exalted about him?” Tibor asked.

“Nothing,” the bird said. “Unless I’m mistaken.”

Tibor cursed softly.

“You’ll have to judge for yourself,” the bird said. “Your luck hasn’t run out.”

“How do you know?”

“That’s one thing I’m certain of.”

“What kind of bird are you?” Tibor asked.

“A blue jay.”

“Are blue jays reliable?”

“Generally,” the bird said.

The jay landed on his shoulder. It spoke gently now. “Who else do you have? I waited years for you. When I heard you singing hymns, I knew. Trust me.”

Tibor nodded. “A hymn-singing bird should be trusted.”

He reminded the bird of the cart’s ruined bearings.

“The slime works,” the bird said. “Come.”

They moved north beneath a clear sky. The world seemed quieter in daylight; fewer creatures emerged. Tibor followed the bird as if it were a star.

Ahead lay hills and a fertile valley. Roads converged. Buildings clustered. A settlement.

“New Brunswick, Idaho,” the bird said.

Tibor’s bearings shrieked again. He hoped for a blacksmith.

The bird explained that nothing was worthless here. Bedsprings became tools. Radios became antifertility devices.

“God,” Tibor said. “They’d rather no children than altered ones?”

“Yes,” the bird said.

As they descended, the bird described grotesque mutations. Tibor grew ill listening. Then the wheel finally tore loose. The cart halted violently.

“That’s it,” Tibor said. Panic surged. He could not move. Could not escape.

The bird offered to fetch help. Tibor wrote a note and watched the jay vanish with it.

He waited.

He prayed.

Nothing answered.

He recited fragments of mass, of requiem, of childhood memory. Still nothing.

Then the sky ignited.

A blazing red disk descended, roaring like molten metal. A face formed—furious, immense.

“You mock me,” it thundered. “Pray!”

“I have no knees,” Tibor said.

The being lifted him. He knelt.

He had legs.

He had arms.

“You are Carleton Lufteufel,” Tibor whispered.

“Pray.”

Tibor began—but argued, explained, justified himself. The God listened impatiently.

“I know,” the God said. “I sent the bird. I caused the worm. I broke your wheel. You have always been in my power.”

Tibor raised his camera and took a photograph.

The image revealed Lufteufel unmistakably.

The God destroyed the photograph.

Then the limbs were taken back.

“Will you end your Pilg?” the God asked.

“No,” Tibor said.

“You must not find my mortal body,” the God warned—and vanished.

Tibor sat alone again, limbless in his cart.

He considered killing the God—but knew he never could have.

The encounter had changed him.

He shouted for help with his bullhorn.

No answer came.

He waited.

